\section{Nine-point QED amplitudes}\label{sec:nine-qed}

Seven photons on one charged line give $7!=5{,}040$ diagrams.
They fit in a single permutation sum. More charged pairs introduce a
different choice: which line emitted each photon? At nine points there
are at most four charged lines. Connecting them gives just an edge, a
three-line path, or the four-line star and path of
Sec.~\ref{sec:eight-four-pairs}. The following formulas keep the
external massive spins open. Labels related by parity and permutations
of equal-helicity photons cover all $128+32+8+2=170$ assignments.

\paragraph{One charged pair.}
Use the ordered word of Eq.~\eqref{eq:eight-ordered-tensor}. Put $c$
at its fourth vertex, $L$ on the positron side and $R$ on the electron
side. Define $B(L;c;R)$ by summing that word over the six orders in
each triple separately. Include the two fermion propagators immediately
adjacent to $G_c$. Thus $B$ is a row spinor times $G_c$ times a column
spinor. The triples can be computed once and reused for different
central photons. If $h$ of the seven photons have negative helicity,
\begin{equation}
 \M_h=-e^7\,\mathcal S_+\mathcal S_-
 \sum_{\substack{c_-=0,1;\ 0\le j\le3\\ 0\le k=h-c_--j\le3}}
 \frac{B(L;c;R)}{j!(3-j)!\,k!(3-k)!}.
 \label{eq:nine-single}
\end{equation}
Here $c_-=1$ selects a negative central photon, $j$ and $k$ count
the negatives on the two sides, and $\mathcal S_\pm$ sums permutations
with unit weight. For each term assign the positive labels in increasing
order to $L,c,R$, then the negative labels in the same way. In Eq.~\eqref{eq:nine-single}, assign a polarization vector to each
external photon before forming the sum. Each permuted leg then carries
that vector into its new position. In particular, the central photon
and the two triples use the same reference assignment throughout the
complete sum. The factorials divide out orders already included in the triples.
For $h=0,1,2,3$ there are respectively $1,3,5,7$ summands. Each labeled
order occurs once because it has a unique fourth vertex. For $h=0$ the
same-helicity formulas in the main text can also be used directly with
seven photons.

\paragraph{Seven positive photons.}
Here the spin dependence factors into
$C=(\sqrt{2}e)^7m\langle\mathbf1\mathbf2\rangle$. Put
$X_i=[i12i]$ and $E_i(L)=\sum_{a\in L}[ia]\langle a1i]$.
For the physical form the ordered correction has only four steps,
\begin{align*}
 R_{89}&=[89]^2\tau_{18}\tau_{19}/\tau_{289},\\
 R_{789}&=\frac{3[78]^2\tau_{17}\tau_{18}}{\tau_{2789}}
 \left(X_9+\frac{\tau_{29}}{\tau_{278}}[X_9-E_9(78)]\right),\\
 R_{6789}&=\frac{6[67]^2\tau_{16}\tau_{17}X_8X_9
  +4 E_6(345)\tau_{26}R_{789}
  -6m^2[67]^2\tau_{26}\tau_{27}R_{89}}{\tau_{1345}},\\
 R_{56789}&=\frac{10[56]^2\tau_{15}\tau_{16} X_7X_8X_9
 +5 E_5(34)\tau_{25}R_{6789}
 -10m^2[56]^2\tau_{25}\tau_{26}R_{789}}{\tau_{134}} .
\end{align*}
For instance, the factor ten chooses the marked pair among the five
spectators of legs 3 and 4. Momentum conservation gives
$\tau_{256789}=\tau_{134}$. With an unnormalized sum over the seven
photon labels, the amplitude is
\begin{equation}
 \M^{\phys}_{+++++++}=\frac{C}{240}\sum_{S_7}
 \frac{[34]^2\,[X_5X_6X_7X_8X_9-m^2R_{56789}]}
 {\displaystyle\prod_{i=3}^9\tau_{1i}\prod_{j=5}^9\tau_{2j}}.
 \label{eq:nine-seven-physical}
\end{equation}
The divisor is $2!5!$. The order inside the terminal pair is immaterial.

For the spurious form put
$Y_i(L)=\sum_{a\in L}\langle3ai]$ and
$y_i=Y_i(1)$. For a photon label $a$, write $\langle3ai]=\langle3a\rangle[ai]$.
The three spectator factors become
\begin{align*}
 Z_9 &= y_9+\tau_{19}Y_9(156)/\tau_{156},\\
 Z_{89} &= y_8y_9-2\tau_{18}Y_8(247)Z_9/\tau_{1569},\\
 Z_{789} &=y_7y_8y_9-3\tau_{17}Y_7(24)Z_{89}/\tau_{2347}.
\end{align*}
The second momentum denominator in the last line satisfies
$\tau_{15689}=\tau_{2347}$. The angle reference now distinguishes
photon 3 from photon 4, giving the divisor $5!$,
\begin{equation}
 \M^{\spur}_{+++++++}=-\frac{C}{120}\sum_{S_7}
 \frac{[34]}{\langle34\rangle\,\tau_{23}\prod_{i=5}^9\tau_{1i}\prod_{i=7}^9\langle3i\rangle}
 \left[\frac{y_5y_6y_7y_8y_9}{\langle35\rangle\langle36\rangle}-
 \frac{10m^2[56]^2Z_{789}}{\tau_{234}}\right].
 \label{eq:nine-seven-spurious}
\end{equation}

A second spurious form fits in a single product per photon permutation.
Choose photon 8 as the reference and 9 as its partner. After the first
three emissions, use momentum conservation to shorten the last two
sandwiches. With the momentum-sum notation of
Eq.~\eqref{eq:eight-ordered}, the result is
\begin{align}
 \M^{\spur}_{+++++++}={}&-C\sum_{S_7}
 \frac{[89]\,\langle813]\,\langle8(13)4]\,\langle8(134)5]\,
                \langle8(279)6]\,\langle8(29)7]}
 {\langle89\rangle\,\tau_{28}\,
  \tau_{13}\tau_{134}\tau_{1345}\tau_{2789}\tau_{289}
  \displaystyle\prod_{i=3}^{7}\langle8i\rangle}.
 \label{eq:nine-seven-product}
\end{align}
The $7\cdot6$ choices of anchor and partner leave $5!$ emission orders.
Compared with Eq.~\eqref{eq:nine-seven-spurious}, this expression is
shorter on the page and carries more explicit momentum insertions.
Pulling the three common angle denominators outside the correction in
Eq.~\eqref{eq:nine-seven-spurious} and resolving the massless insertions
reduces the HEPCAT summand score from 372 to 208.

\paragraph{Six positive photons and one negative.}
There is a shorter spurious form than the three middle-vertex terms in
Eq.~\eqref{eq:nine-single}. Insert $h=9$ into the seven gaps of an
ordered word $(i_1,\ldots,i_6)$ of positive photons. As in
Sec.~\ref{sec:eight-one-minus}, use $h$ as their common reference and
choose $r=i_3$ for the negative photon. Write
\[
 Q_j=p_1+\sum_{a=1}^{j}p_{i_a},\qquad
 U=\langle\mathbf1\mathbf2\rangle,\qquad
 H=\langle\mathbf1h\rangle\langle h\mathbf2\rangle .
\]
Since $[i|p_h|h\rangle=0$, all seven gaps give the same orbital factors.
Before the insertion, the propagators accumulate from leg 1. After the
insertion, count them back from leg 2. This gives one product times a
sum over the position of the negative photon,
\begin{align}
 \M^{\spur}_{++++++-}
  =\frac{(\sqrt{2}e)^7}{m}\sum_{S_6}
  \left(\prod_{a=1}^{6}\frac{[i_a|Q_{a-1}|h\rangle}
                           {\langle h i_a\rangle}\right)
  \sum_{j=0}^{6}
  \frac{H+U\,[r|Q_j|h\rangle/[rh]}
       {\displaystyle\prod_{a=1}^{j}\tau_{1 i_1\cdots i_a}
        \prod_{a=j+1}^{6}\tau_{2 i_a\cdots i_6}} .
 \label{eq:nine-one-gaps}
\end{align}
For example the gap after $i_2$ has left denominators $\tau_{1i_1}$
and $\tau_{1i_1i_2}$ and four right denominators ending with
$\tau_{2i_6}$. All spurious denominators in this form are brackets
between external photons.

\paragraph{Two or more charged pairs.}
For two pairs use $(1,2),(4,3)$ in electron--positron order. With
three or four pairs use $(1,2),(3,4),\ldots$. Let $\mathcal A_e$ and
$\mathcal A_E$ denote the two unnormalized alternating sums over the
electron and positron labels. For a tree $t$ and sets $L_v$ of real
photons on its charged lines, write
\begin{equation}
 K_t(L_1|\cdots|L_p)=(-1)^p e^7
 \left(\prod_{vw\in t}\frac{g_{\mu_{vw}\nu_{vw}}}{q_{vw}^2}\right)
 \prod_{v=1}^p T_v(L_v,\text{exchanges incident on }v).
 \label{eq:nine-tree}
\end{equation}
Use the complete open tensor of Eq.~\eqref{eq:eight-ordered-tensor}
on each line. Each exchanged photon carries the sum of the external
momenta on its far side. Its two indices in Eq.~\eqref{eq:nine-tree}
belong to its two incident tensors. The common sign follows from the
seven vertices $-ie\gamma^\mu$, the six propagators (including $p-1$
photon propagators) and the overall $i$ removed from the amplitude.
Permuting the external fermions supplies the sign of each matching.

In Table~\ref{tab:nine-occupations} a string $L|M|R$ stands for
$K_{\rm path}(L|M|R)$, and a zero denotes an empty set. Use successive
letters $a,b,\ldots$ for positive photons and $h,i$ for negatives.
The amplitude is the sum of the indicated terms acted on by
$\mathcal A_e\mathcal A_E\mathcal S_+\mathcal S_-$.
For example, the denominator 120 in the first row removes the $5!$
orders of the five identical-helicity labels already summed inside
one tensor. The denominator 12 in $abc|de$ removes $3!2!$ redundant
permutations. No factor for exchange of the two lines survives in
this example since their photon occupancies differ.
\begin{table}[htbp]
\centering\small
\begin{tabular}{cl}
\toprule
Helicities & Terms and their divisors \\
\midrule
$+++++$ & $(abcde|0)/120+(abcd|e)/24+(abc|de)/12$\\
$++++-$ & $[(abcdh|0)+(abcd|h)]/24+[(abch|d)+(abc|dh)]/6+(abh|cd)/4$\\
$+++--$ & $[(abchi|0)+(abc|hi)]/12+(abch|i)/6$\\
 & $\quad+[(abhi|c)+(ab|chi)]/4+(abh|ci)/2$\\
\midrule
$+++$ & $(abc|0|0)/6+[(ab|c|0)+(ab|0|c)+(a|bc|0)+(a|b|c)]/2$\\
 & $\quad +(0|abc|0)/12$\\
$++-$ & $[(abh|0|0)+(ab|h|0)+(ab|0|h)+(a|h|b)+(h|ab|0)]/2$\\
 & $\quad +(ah|b|0)+(a|bh|0)+(a|b|h)+(ah|0|b)+(0|abh|0)/4$\\
\bottomrule
\end{tabular}
\caption{Photon allocations on two charged lines (first three rows)
and on a three-line path (last two rows).}\label{tab:nine-occupations}
\end{table}

For four pairs and one photon $a$, label the center of the star first
and list the four vertices of the path in order. The alternating sums
act on four terms,
\begin{align}
 \M={}&\mathcal A_e\mathcal A_E\left[
 \frac{K_{\rm star}(a|0|0|0)}6+\frac{K_{\rm star}(0|a|0|0)}2
 +K_{\rm path}(a|0|0|0)+K_{\rm path}(0|a|0|0)\right].
 \label{eq:nine-four-pairs}
\end{align}
In the first star all three leaves are equivalent. In the second only
the two empty leaves can be exchanged. A marked vertex breaks the
reflection symmetry of the path. These four possibilities account for
17,280 distinct ordered tree diagrams.

\paragraph{Pole choices and checks.}
For the physical representation use the massive electron momentum of
the emitting line as the reference of each photon on that line. This
gives reference denominators proportional to $2p_e\cdot p_i$, a
physical electron virtuality. For a representation with spurious poles,
take a photon of opposite helicity when possible, and otherwise use
the next photon cyclically. Apply the reference prescription to the
assigned labels. When the process has only one real photon $k$, emitted by line $(a,b)$, the spurious prescription
uses $Q=p_a+p_b$. Take $Q|k]$ as the angle reference, or its parity
conjugate for negative helicity. The new divisor is
$2k\cdot Q=s_{abk}-s_{ab}$. At finite propagator virtualities the
residue on this divisor vanishes by the Ward identity for that line. Each line sums all orders
of its attachments; the Ward identity follows by telescoping adjacent
fermion propagators and using the endpoint Dirac equations.

A terminal line emitting one real photon admits a shorter tensor. Write
$G=\gamma^\mu$ for its exchanged attachment and $D=\not k\not\epsilon$
for the real photon. Anticommuting the endpoint momenta through the
emission vertex in the two orders gives
\begin{align}
 T_{ab}^{\mu}(k)={}&
 \left(\frac{2p_a\cdot\epsilon}{\tau_{ak}}
       -\frac{2p_b\cdot\epsilon}{\tau_{bk}}\right)\bar v_b G u_a
 +\bar v_b\left(\frac{GD}{\tau_{ak}}+\frac{DG}{\tau_{bk}}\right)u_a .
 \label{eq:nine-leaf}
\end{align}
Both terms vanish separately for $\epsilon=k$. In the physical gauge
$p_a\cdot\epsilon=0$. In the single-photon spurious gauge
$p_b\cdot\epsilon=-p_a\cdot\epsilon$. Thus both gauges remove one
independent orbital contraction. An unmarked leaf contributes just
$\bar v_b\gamma^\mu u_a$.

For one through four pairs the numbers of distinct physical channels
are 126, 188, 186 and 172. The ordered diagram counts are 5,040, 13,440,
18,144 and 17,280. These counts concern labeled diagrams; the number
of summands exhibited above measures the effect of the permutation sums.

We checked the printed occupation divisors directly by expanding their
permutations. The resulting attachment words and particle matchings
agree with the independently generated diagram lists. In addition, for
each of the 20 helicity counts, the two printed tensor prescriptions
agreed with an independent-gauge current at 32 coherent-spin points.
The maximum relative discrepancies in the signed amplitudes were
$7.64\times10^{-13}$ and $2.83\times10^{-12}$, respectively. At four
separate points per helicity count the corresponding polarized rates, summed
over every massive spin, agreed with MadGraph within $1.41\times10^{-13}$
and $1.11\times10^{-12}$. We also tested each of the 170 labeled photon
helicity assignments directly at two new points, with maximum discrepancies
of $3.81\times10^{-13}$ and $5.13\times10^{-12}$. These samples used
$m=1$ and $e=1/3$.
